% ====================================================================
% method.tex (dissertation example)
% ====================================================================

\section{Method}
\label{sec:method}

\subsection{Approach}

The questions were investigated by simulation and not on hardware. On a real
microcontroller the input noise cannot be set to a chosen value, and it is
mixed with the converter's other imperfections, so the effect of the noise
alone could not be separated out. A simulation allows the noise to be set
exactly, including to zero, and the true input is known at every sample, so
the error can be calculated directly. The cost of this choice is that the
results describe an ideal converter, which is considered in
\cref{sec:discussion}.

\subsection{The model}

The model has three parts, which follow \cref{fig:oversamplingChain}. The
input is a voltage expressed in units of one LSB. Gaussian noise with a chosen
rms value is added to it, and the sum is rounded to the nearest whole code and
limited to the range of a 12-bit converter. \Cref{lst:adcModel} shows this
part. Each output sample is then the mean of $M$ successive codes.

% ====================================================================
% adcModel.tex
% --------------------------------------------------------------------
% The converter model, shown as an excerpt. The line range picks the
% function out of the real source file, so the listing always shows the
% code that produced the results. Update the range if the file changes.
% ====================================================================

\lstinputlisting[
  language=Python,
  caption={The model of an ideal converter with noise at its input.},
  label={lst:adcModel},
  linerange={26-30},
  firstnumber=26
]{assets/code/python/oversampling_sim.py}


The model was written in Python using only its standard library, so that it
can be run without installing anything further, and the random number
generator is given a fixed seed so that the results repeat exactly. The full
source is listed in Appendix~\ref{app:source}. The parameters are in
\cref{tab:parameters}.

\begin{table}[htbp]
\centering
\begin{tabularx}{0.9\textwidth}{@{} X l @{}}
\toprule
\textbf{Parameter} & \textbf{Value} \\
\midrule
Converter resolution, $N$                 & 12 bits \\
Oversampling ratios, $M$                  & 1, 4, 16, 64, 256 \\
Input noise, $\sigma_n$                   & 0, 0.25, 0.5 and \qty{1.0}{LSB} rms \\
Sine test: amplitude                      & \qty{2000}{LSB}, centred on mid-scale \\
Sine test: frequency                      & \qty{3.71}{\percent} of the output sample rate \\
Sine test: output samples per run         & 4096 \\
Static test: input levels per run         & 2000, placed at random between two codes \\
\bottomrule
\end{tabularx}
\caption{Parameters of the simulation.}
\label{tab:parameters}
\end{table}

\subsection{Measuring the resolution}

In both tests the error of each output sample is the difference between the
averaged codes and the average of the true input over the same samples. The
rms value of that error over the run, $e_{\mathrm{out}}$, is converted to an
effective number of bits by comparing it with \cref{eq:quantisationNoise},

\begin{equation}
  \label{eq:enobMeasured}
  \mathrm{ENOB} = N - \log_2\!\left(\frac{e_{\mathrm{out}}\sqrt{12}}{q}\right),
\end{equation}
so that an error equal to that of an ideal 12-bit converter gives exactly 12
bits.

\subsection{Test 1: sine wave input}

The first test uses a sine wave covering most of the input range, with
\qty{0.5}{LSB} rms of noise, and measures the resolution at each oversampling
ratio. This is the case that \cref{eq:enobTheory} describes, so it serves as
the check on the model as well as the first result. The frequency was chosen
so that the run does not contain a whole number of cycles, which avoids the
same few codes being exercised repeatedly.

\subsection{Test 2: static inputs}

The second test holds the input steady for the whole of each average, which
is the condition under which the method is expected to fail without noise. A
single level would not be representative, as an input that happens to sit at
the centre of a step has no error to begin with, so each run uses 2000 levels
spread at random across one step and takes the rms error over all of them.
The test was repeated at each of the four noise levels in
\cref{tab:parameters}.
